\chapter{Yannakakis: cadena}
\label{cap:yannakakis-chain}

\begin{intuicion}
Yannakakis (1981) tiene dos partes: \emph{semi-join reducer} y
\emph{enumerator}. El reducer baja por el join tree (hojas $\to$
raíz) y sube (raíz $\to$ hojas), eliminando en cada paso las
tuplas que no participan de ningún match completo. El enumerator
después recorre las tuplas que sobrevivieron y enumera el join
final ---ya sin desperdicio. Para hipótesis con shape \emph{chain},
el join tree es lineal, las dos pasadas son simples filtros, y el
algoritmo termina en \(O(\)input + output\()\). Este capítulo lo
desarrolla con un ejemplo concreto y prueba la correctitud.
\end{intuicion}

\section{Setup: \emph{bags} y la chain}

Antes de Yannakakis, el motor materializa las relaciones de la
hipótesis como \emph{bags}. Para una chain
\((L_0^V \to_{L_0^E} L_1^V \to_{L_1^E} \dots \to_{L_{n-1}^E} L_n^V)\),
\(\mathrm{bag}_i\) contiene las aristas:
\[
B_i = \{(u, v, \tau) : (u, v, \tau) \in \E,\;
                     \mathrm{type}(u) = L_i^V,\;
                     \mathrm{type}_E(\cdot) = L_i^E,\;
                     \mathrm{type}(v) = L_{i+1}^V\}.
\]

\noindent
Cada \(B_i\) es la lista de aristas que podrían llenar el paso
\(i\). El join chain es:
\[
B_0 \bowtie_{B_0.\mathrm{dest} = B_1.\mathrm{origin}} B_1
   \bowtie_{B_1.\mathrm{dest} = B_2.\mathrm{origin}} \dots
   \bowtie_{B_{n-2}.\mathrm{dest} = B_{n-1}.\mathrm{origin}} B_{n-1}.
\]

\section{La semi-join reducer}

La idea: filtrar cada \(B_i\) hasta que sólo queden aristas que
participan de \emph{algún} match completo.

\subsection{Pasada bottom-up}

\begin{algorithm}[H]
\caption{Pasada 1 (bottom-up): de \(B_{n-1}\) hacia \(B_0\)}
\For{\(i\) from \(n - 2\) \textbf{down to} \(0\)}{
    \(\texttt{live\_origins} \leftarrow \{r.\mathrm{origin} : r \in B_{i+1}\}\)\;
    \(B_i \leftarrow \{r \in B_i : r.\mathrm{dest} \in \texttt{live\_origins}\}\)\;
}
\end{algorithm}

\noindent
Después de esta pasada: cada arista de \(B_i\) tiene su
\texttt{dest} matched con al menos una arista de \(B_{i+1}\). Es
decir, hay una continuación posible hacia adelante.

\subsection{Pasada top-down}

\begin{algorithm}[H]
\caption{Pasada 2 (top-down): de \(B_0\) hacia \(B_{n-1}\)}
\For{\(i\) from \(1\) \textbf{to} \(n - 1\)}{
    \(\texttt{live\_dests} \leftarrow \{r.\mathrm{dest} : r \in B_{i-1}\}\)\;
    \(B_i \leftarrow \{r \in B_i : r.\mathrm{origin} \in \texttt{live\_dests}\}\)\;
}
\end{algorithm}

\noindent
Después de esta pasada: cada arista de \(B_i\) tiene su
\texttt{origin} matched con al menos una arista de \(B_{i-1}\) que
sobrevivió. Hay una continuación posible hacia atrás.

\subsection{Diagrama}

\begin{figure}[H]
\centering
\begin{tikzpicture}[
    node distance=0.3cm,
    every node/.style={font=\scriptsize, align=center},
    bag/.style={draw=accent, thick, rounded corners=1mm,
                minimum width=2.4cm, minimum height=1cm,
                fill=accentlight}
]
\node[bag] (b0) {$B_0$\\\itshape User--Auth--Host};
\node[bag, right=of b0] (b1) {$B_1$\\\itshape Host--Connect--IP};
\node[bag, right=of b1] (b2) {$B_2$\\\itshape IP--Beacon--Domain};

\draw[->, >=Latex, accent, thick, dashed] (b1.south) to[out=-90,in=-90, looseness=1.2] node[below, yshift=-0.2cm]{semi-join} (b0.south);
\draw[->, >=Latex, accent, thick, dashed] (b2.south) to[out=-90,in=-90, looseness=1.2] (b1.south);
\draw[->, >=Latex, accent!60, thick] (b0.north) to[out=90,in=90, looseness=1.2] node[above, yshift=0.2cm]{semi-join} (b1.north);
\draw[->, >=Latex, accent!60, thick] (b1.north) to[out=90,in=90, looseness=1.2] (b2.north);
\end{tikzpicture}
\caption{Las dos pasadas del semi-join reducer. Bottom-up
(flechas dashed abajo) elimina rows cuyo \texttt{dest} no aparece
como \texttt{origin} en el siguiente bag; top-down (flechas
arriba) hace lo simétrico. Después, toda row sobreviviente
participa de al menos un full join.}
\label{fig:semijoin-passes}
\end{figure}

\section{Enumerar lo que quedó}

\begin{algorithm}[H]
\caption{Yannakakis chain matcher}
\KwIn{\(B_0, B_1, \dots, B_{n-1}\) (bags reducidos)}
\KwOut{Vec<HuntResult>}
\For{cada \texttt{row} en \(B_0\)}{
    \texttt{walk}(\texttt{row.origin}, \texttt{row.dest}, \texttt{row.timestamp}, \(1\))\;
}
\BlankLine
\Comment{\texttt{walk(prev\_dest, curr\_dest, t\_first, step)}}
\If{\(\mathrm{step} = n\)}{
    \texttt{emit}(path acumulado)\;
    \Return\;
}
\For{cada \(r \in B_{\mathrm{step}}\) con \(r.\mathrm{origin} = \mathrm{prev\_dest}\)}{
    chequear ventana temporal y monotonicidad\;
    \texttt{walk}(\(r.\mathrm{origin}, r.\mathrm{dest}, t_{\mathrm{first}}, \mathrm{step} + 1\))\;
}
\end{algorithm}

\noindent
La enumeración es esencialmente un DFS, \emph{pero} sobre los bags
ya reducidos. Cada arista que el walk considera está garantizada
de participar en algún match (por las dos pasadas previas), así que
el walk \emph{nunca} explora ramas muertas.

\section{Correctitud}

\begin{theorem}[Yannakakis chain correcto]
Después de las dos pasadas semi-join, el conjunto de matches
enumerados por \texttt{walk} es exactamente el join completo
\(B_0 \bowtie B_1 \bowtie \dots \bowtie B_{n-1}\).
\end{theorem}

\begin{proof}[Idea]
Soundness: cada \texttt{emit} corresponde a una secuencia de
aristas \((r_0, r_1, \dots, r_{n-1})\) con \(r_i.\mathrm{dest} =
r_{i+1}.\mathrm{origin}\) (chequeado en el walk) y cada \(r_i \in
B_i\) (porque iteramos sobre los bags). Eso es exactamente la
definición de \(\bowtie\).

Completeness: cualquier match \((r_0, \dots, r_{n-1})\) del join
original tiene cada \(r_i\) que \emph{sobrevivió} la pasada bottom-up
(porque tenía continuación a \(r_{i+1}\)) y la top-down (porque
había \(r_{i-1}\)). Por lo tanto cada \(r_i\) sigue en \(B_i\)
después del reducer; \texttt{walk} los enumera. \(\square\)
\end{proof}

\section{Análisis de complejidad}

\subsection{Costo del semi-join reducer}

Cada pasada itera todas las aristas de cada bag una vez:
\(O(\sum_i |B_i|) = O(|\E|_{\text{relevantes}})\). El conjunto
\texttt{live\_origins} (resp. \texttt{live\_dests}) se construye
con un \texttt{HashSet} en una pasada lineal, y los chequeos
\texttt{r.dest \(\in\) live\_origins} son \(O(1)\).

Total: \(O(|\E|_{\text{relevantes}})\) por pasada, dos pasadas
\(\Rightarrow\) \(O(|\E|)\) total.

\subsection{Costo del enumerator}

Cada path completo (cada match) toma \(O(n)\) trabajo (push +
emit). En total \(O(n \cdot |Q_{\text{out}}|)\). Para hipótesis
chain con \(n\) chico (\(\leq 6\)), \(O(|Q_{\text{out}}|)\).

\subsection{Total}

\(O(|\E| + |Q_{\text{out}}|)\). Lineal en input + output. Es la
promesa de Yannakakis para \(\alpha\)-acíclicas.

\section{Comparación con DFS}

DFS para la misma hipótesis:

\begin{itemize}
    \item Itera $|Q_{\text{out}}|$ paths exitosos.
    \item Itera además todos los paths fallidos (que el reducer
    eliminó).
\end{itemize}

\noindent
En grafos donde la mayoría de los paths parciales no se completan
(filtro temporal, etiquetas selectivas), el DFS gasta mucho en
exploración inútil. Yannakakis lo evita haciendo el chequeo de
``hay continuación'' como semi-join lineal antes de entrar al
walk.

\subsection{Ejemplo concreto}

Hipótesis chain de 3, \(|B_0| = 10^4\), \(|B_1| = 10^5\), \(|B_2|
= 10^6\), \(|Q_{\text{out}}| = 10^3\) (matches reales).

\begin{itemize}
    \item DFS: explora \(\sim 10^4 \cdot 10 \cdot 10 = 10^6\) paths
    parciales (asumiendo grado promedio 10 en cada paso).
    \item Yannakakis: 2 pasadas \(O(|\E|_{\text{relevantes}}) =
    O(10^6)\). Después enumera \(10^3\) paths.
\end{itemize}

\noindent
Mismo orden de magnitud en CPU total, pero Yannakakis paga el costo
en operaciones simples (semi-join lineal con HashSet) en vez de
recursión profunda con \emph{branch mispredictions} y
\emph{cache misses}. Bench típico: \(2\times\) más rápido para
hipótesis con muchos paths fallidos.

\section{Cuando Yannakakis pierde}

\begin{itemize}
    \item \textbf{Hipótesis con cap saturado}: si el motor produce
    \(\geq 10\,000\) results, el costo del enumerator domina; el
    semi-join reducer es overhead constante. DFS con cap-aware
    early-exit es comparable.
    \item \textbf{Hipótesis con muy pocos starts}: si NLF + k-hop
    ya filtraron hasta dejar 10 starts válidos, el costo del DFS
    desde 10 starts es trivial; el semi-join sobre todo el grafo
    es caro innecesariamente.
\end{itemize}

\noindent
El \emph{planner} (capítulo~\ref{cap:planner}) usa estos
heurísticos para decidir cuándo dispatchar Yannakakis vs DFS.

\section{Resumen}

\begin{itemize}
    \item Yannakakis chain: dos pasadas semi-join (bottom-up +
    top-down) + enumerator.
    \item Después del reducer, cada arista en cada bag participa
    de algún match completo. El enumerator no explora ramas muertas.
    \item Complejidad: \(O(|\E| + |Q_{\text{out}}|)\). Lineal en
    input + output.
    \item Gana sobre DFS cuando hay muchos paths parciales que se
    abortan; pierde cuando NLF + k-hop ya filtraron al 90\%.
\end{itemize}

\begin{ejercicios}
\begin{enumerate}
    \item \starred{} Para la hipótesis
    \texttt{User -[Auth]-> Host -[Connect]-> IP} con \(|B_0| =
    1000\), \(|B_1| = 5000\), y los matches finales son 50, calculá
    cuántas filas sobreviven a cada pasada del reducer en el peor
    caso.

    \item Implementá la pasada bottom-up con un \texttt{HashSet}
    auxiliar. ¿Cuál es la complejidad de espacio?

    \item \starred{} Demostrá: si después de la pasada bottom-up
    \(B_0\) queda vacío, entonces \(Q\) está vacío (no hay matches).
    ¿Es \emph{cierto el converse}?

    \item Diseñá un ejemplo donde la pasada top-down NO elimina
    ninguna fila adicional (es no-op). ¿Cuándo ocurre?

    \item Compará Yannakakis chain con DFS bien-podado para la
    hipótesis del ejercicio 1. ¿Cuántos paths visita cada uno?
    ¿Bajo qué condiciones gana cada uno?

    \item Si la hipótesis tiene un predicado complejo en uno de
    los pasos (\texttt{HAVING count(distinct)} \(>\) 5), ¿cómo lo
    integrás en Yannakakis chain? ¿Antes, durante, o después del
    reducer?
\end{enumerate}
\end{ejercicios}
